\chapter*{Introduction Générale}
\addcontentsline{toc}{chapter}{Introduction Générale}

Depuis la nuit des temps, l'investissement financier est guidé par une quête universelle : maximiser le gain tout en minimisant la probabilité de perte. Cependant, pendant des siècles, cette gestion a reposé sur l'intuition, l'expérience empirique et parfois la spéculation pure. Ce n'est qu'au milieu du XXe siècle que la gestion de portefeuille est entrée dans l'ère scientifique avec les travaux fondateurs de Harry Markowitz en 1952. En formalisant mathématiquement le compromis entre le rendement espéré et le risque (mesuré par la variance), Markowitz a jeté les bases de la \textit{Théorie Moderne du Portefeuille} (MPT). Pour la première fois, la diversification n'était plus un simple conseil de bon sens, mais une règle quantitative rigoureuse démontrant que le risque d'un portefeuille dépend de la covariance des actifs qui le composent plutôt que de la simple somme de leurs risques individuels.

Pourtant, malgré son élégance mathématique et son adoption massive par l'industrie financière, le modèle de Markowitz repose sur des hypothèses fortes et souvent irréalistes dans la pratique. Il suppose notamment que les rendements des actifs suivent une loi normale, que les investisseurs sont strictement rationnels et que la structure statistique des rendements passés (moyennes, variances et covariances) restera stable dans le futur. Les crises financières successives de ce début de siècle (la bulle Internet en 2001, la crise des Subprimes en 2008 et le krach lié à la pandémie de COVID-19 en 2020) ont mis en lumière les limites critiques de ces modèles classiques. Sur les marchés réels, les rendements ne sont pas gaussiens ; ils affichent des « queues épaisses » (fat tails) qui augmentent la fréquence des événements extrêmes (les fameux « Cygnes Noirs » de Nassim Nicholas Taleb). De plus, l'optimiseur classique de Markowitz s'est révélé être un « amplificateur d'erreurs » : de légères modifications dans les données historiques d'entrée conduisent à des allocations de portefeuille extrêmes et impraticables.

Parallèlement, les marchés financiers ont connu une mutation technologique sans précédent. L'ère de la numérisation a entraîné une explosion du volume, de la variété et de la vitesse de propagation des données financières (le Big Data). Face à cette masse de données complexes, changeantes et souvent non-linéaires, l'analyse humaine et les outils statistiques classiques s'avèrent limités. C'est dans ce contexte que l'Intelligence Artificielle (IA) et plus spécifiquement le Machine Learning (Apprentissage Automatique) et le Deep Learning (Apprentissage Profond) se sont imposés comme des outils de rupture pour la finance quantitative.

Les réseaux de neurones artificiels, et en particulier les architectures récurrentes de type \textbf{LSTM (Long Short-Term Memory)}, ont été spécialement conçus pour modéliser des séries temporelles. Contrairement aux réseaux de neurones traditionnels, les LSTM possèdent une mémoire interne capable de capturer des dépendances temporelles à long terme et d'assimiler les dynamiques non-linéaires des marchés sans imposer d'hypothèse statistique rigide \textit{a priori} sur la distribution des prix. 

L'objet de ce Projet de Fin d'Études est de concevoir, d'implémenter et de tester une méthodologie hybride d'allocation d'actifs. Cette approche combine la puissance prédictive de l'apprentissage profond (modèle LSTM) avec le cadre d'optimisation robuste de la Théorie Moderne du Portefeuille. L'idée centrale est de substituer les rendements historiques statiques (habituellement utilisés pour estimer le vecteur des rendements espérés dans l'optimiseur de Markowitz) par des prédictions dynamiques générées par le réseau LSTM.

Pour répondre à cette problématique et évaluer l'efficacité de cette démarche, notre étude s'articulera autour de quatre chapitres principaux :

\begin{itemize}
    \item \textbf{Le Chapitre 1} posera le cadre conceptuel en présentant la structure des marchés financiers, les instruments utilisés, la modélisation mathématique classique du couple rendement-risque, ainsi que l'hypothèse d'efficience des marchés et les limites de la finance traditionnelle face au risque.
    \item \textbf{Le Chapitre 2} détaillera la transition de la théorie classique vers l'intelligence artificielle. Nous y étudierons en profondeur le modèle de Markowitz, sa formulation mathématique, la dérivation de la frontière efficiente, le MEDAF/CAPM, ainsi que les critiques et limites empiriques de ces approches classiques.
    \item \textbf{Le Chapitre 3} introduira les fondements de l'Intelligence Artificielle en finance quantitative. Nous comparerons les différents algorithmes de prédiction (Régression linéaire, ARIMA, Random Forest) et détaillerons l'architecture mathématique des réseaux de neurones récurrents LSTM. Enfin, nous décrirons le fonctionnement de notre stratégie hybride LSTM + Markowitz.
    \item \textbf{Le Chapitre 4} sera consacré à l'étude empirique. Nous y présenterons la collecte de données sur 10 ans (2015-2024) pour un portefeuille composé de grandes capitalisations américaines (AAPL, MSFT, GOOGL, JPM, PG) via l'API \texttt{yfinance}. Nous exposerons ensuite l'architecture du modèle LSTM implémenté, les résultats de son entraînement, le tracé de la frontière efficiente et, enfin, le backtesting comparatif des performances de notre stratégie face au modèle classique et à l'indice S\&P 500.
\end{itemize}

À travers ce travail, nous tenterons de déterminer si l'intégration de modèles prédictifs d'intelligence artificielle permet d'améliorer significativement le couple rendement-risque des portefeuilles réels et de surperformer les méthodes traditionnelles d'allocation d'actifs.
